\documentclass[11pt, a4paper]{article}
\usepackage[margin=2.5cm]{geometry}
\usepackage{amsmath, amssymb, amsthm}
\usepackage{fontenc}
\usepackage{inputenc}
\usepackage{hyperref}
\usepackage{xcolor}
\usepackage{enumitem}
\usepackage{titlesec}
\usepackage{booktabs}
\usepackage{graphicx}
\usepackage{microtype}

\hypersetup{
  colorlinks=true,
  linkcolor=blue!60!black,
  urlcolor=blue!70!black,
}

\definecolor{cubblue}{RGB}{68, 136, 255}
\definecolor{cuborange}{RGB}{255, 170, 68}
\definecolor{cubgreen}{RGB}{68, 255, 136}
\definecolor{cubpink}{RGB}{255, 68, 136}

\title{
  \vspace{-1cm}
  \textbf{Cubical Type Theory Visualizer} \\
  \large Design Improvements \& Conceptual Clarity Roadmap
}
\author{Oguri Cap \\ \small cubical-viz project}
\date{March 2026}

\begin{document}
\maketitle

\begin{abstract}
This document proposes a set of improvements to the \texttt{cubical-viz} interactive
visualizer for cubical type theory. We identify conceptual gaps in the current
visualizations, suggest richer visual encodings, and outline new features that would
make the tool substantially more useful for learning and research. The proposals are
organized from quick wins to longer-term ideas.
\end{abstract}

\tableofcontents
\vspace{1em}

% ────────────────────────────────────────────────────────────────────────────
\section{Current State}
% ────────────────────────────────────────────────────────────────────────────

The visualizer currently shows five concepts:
\begin{enumerate}
  \item \textbf{Path Types} — A path as a function $\mathbb{I} \to A$.
  \item \textbf{Composition} — Filling a square from three faces.
  \item \textbf{Transport} — Coercing an element along a path in a type family.
  \item \textbf{Kan Filling} — Constructing the missing lid of an open box.
  \item \textbf{Glue Types} — Gluing types via equivalences to prove univalence.
\end{enumerate}

Each rule has a KaTeX-rendered typing judgment, a prose description, a Three.js 3D
animation, and bidirectional hover highlighting between judgment terms and 3D
objects. The visual style is flat (\texttt{MeshBasicMaterial}), clean, and dark-themed.

% ────────────────────────────────────────────────────────────────────────────
\section{Conceptual Gaps to Address}
% ────────────────────────────────────────────────────────────────────────────

\subsection{The Interval $\mathbb{I}$ Is Under-Explained}

The interval is the foundation of everything, but its special status as a
\emph{non-inductive} primitive with face lattice $\{0, 1, i, 1-i, \ldots\}$ is not
conveyed. The current interval visualization is a static line with endpoints $i_0$
and $i_1$. This misses several key ideas:

\begin{itemize}
  \item The \emph{de Morgan algebra} structure: $i \land j$, $i \lor j$, $1-i$.
  \item The distinction between \emph{internal} (the interval object $\mathbb{I}$)
    and \emph{external} (the meta-level of terms).
  \item Face maps: substituting $0$ or $1$ for $i$ corresponds to selecting
    endpoints; this should be visually interactive (drag a point along the line and
    watch the fiber change).
\end{itemize}

\paragraph{Proposed fix.} Animate a point sliding along the interval and show a
\emph{fiber} — a copy of the type $A(i)$ — changing as $i$ moves. Use color to
encode the value of $i \in [0,1]$.

\subsection{Composition Needs a ``Box'' Not Just an Arc}

Currently, Composition is shown as an arc animation. The key intuition is that
composition \emph{fills a square}: given three sides (two paths and a connecting
path at $i=0$), we get the fourth side. The 3D scene should show a literal square
with three sides drawn and one side being animated in.

\begin{itemize}
  \item Draw a unit square $[0,1]^2$ in 3D space.
  \item Color the three \emph{given} edges in distinct colors corresponding to
    $u_0$, $u_1$, and $p$ (the base path).
  \item Animate the fourth edge (the composition) being filled from $i=0$ to $i=1$.
  \item The formula $\mathrm{comp}^i\,A\,[i=0 \mapsto u_0, i=1 \mapsto u_1]\,p$
    should map directly onto the visual.
\end{itemize}

\subsection{Transport Should Show the Fiber Changing}

Transport moves an element along a path in a \emph{type family}. The current
visualization shows two spheres (for $A(i_0)$ and $A(i_1)$) with a path arc
between them. This is correct but loses the key idea that $A$ is a
\emph{varying} family: the type itself changes as $i$ moves.

\begin{itemize}
  \item Show the path $A : \mathbb{I} \to \mathcal{U}$ as a \emph{sequence of
    shapes morphing} across the screen (e.g., a sphere deforming into a torus).
  \item Show the element $a : A(i_0)$ sitting inside the left fiber, and
    $\mathrm{transp}^i\,A\,a : A(i_1)$ appearing in the right fiber.
  \item Animate the element \emph{moving through the deforming fibers} as $i$
    increases from $0$ to $1$.
\end{itemize}

\subsection{Kan Filling: Clarify ``Open Box'' vs.\ ``Filled Box''}

The current animation fades in the top face. The term \emph{Kan filling} comes from
Kan complexes in simplicial set theory, and the key visual is an \emph{open box}
(five faces of a cube) with the \emph{missing lid} that gets filled.

\begin{itemize}
  \item The box should start visibly \emph{open} at the top --- not just with an
    invisible face but with clearly missing edges.
  \item The \emph{face formula} $\varphi$ should be visualized: $\varphi$ encodes
    \emph{which} faces are given. Draw $\varphi$ as a highlighted subset of faces.
  \item Animate the lid ``snapping in'' or being ``extruded'' from the partial data,
    not just fading in.
  \item The partial element $u$ (the system of given faces) should be individually
    highlighted and hoverable.
\end{itemize}

\subsection{Glue Types: Equivalence Is Not a Path}

The current GlueType visualization draws an arc from type $A$ (a sphere) to type
$B$ (a torus) and animates a point sliding along it. This looks like \emph{transport
along a path}, not an \emph{equivalence}. An equivalence $e : T \simeq A$ is a pair
of maps with homotopy inverses, which is fundamentally different.

\begin{itemize}
  \item Replace the arc with \emph{two arrows} going in both directions ($f : T \to
    A$ and $g : A \to T$) to emphasize the bijectivity.
  \item Show a point in $A$ mapped to $B$ and back, ending at the same point
    (homotopy inverse visualization).
  \item The \emph{Glue} type should be shown as a \emph{blend} or \emph{overlap
    region} between $A$ and $T$, not just a path connecting them.
  \item Add an annotation for the face formula $\varphi$: when $\varphi = 1$,
    we are entirely in $T$; when $\varphi = 0$, we are in $A$.
\end{itemize}

% ────────────────────────────────────────────────────────────────────────────
\section{Visual Design Improvements}
% ────────────────────────────────────────────────────────────────────────────

\subsection{Color Consistency Across Rules}

Currently each rule invents its own color palette. A consistent scheme would help
the user build intuition across pages.

\begin{center}
\begin{tabular}{lll}
\toprule
\textbf{Concept} & \textbf{Suggested Color} & \textbf{Hex} \\
\midrule
Interval / time dimension & Cyan & \texttt{\#44DDFF} \\
Type family $A$ & Blue & \texttt{\#4488FF} \\
Starting element / base point & Yellow & \texttt{\#FFDD44} \\
Result / transported element & Orange & \texttt{\#FF8844} \\
Partial data / face formula & Green & \texttt{\#44FF88} \\
Composition / fill result & Pink & \texttt{\#FF44AA} \\
Equivalence arrow & Purple & \texttt{\#AA44FF} \\
\bottomrule
\end{tabular}
\end{center}

\subsection{Dimension Labels on Axes}

The 3D scene has no axis labels. Users cannot tell which axis corresponds to which
dimension variable ($i$, $j$, $k$). Add thin labeled axes:
\begin{itemize}
  \item $i$-axis: horizontal (red), labeled ``$i : \mathbb{I}$''.
  \item $j$-axis: vertical (green), labeled ``$j : \mathbb{I}$''.
  \item $k$-axis: depth (blue), labeled ``$k : \mathbb{I}$''.
\end{itemize}
These should be toggleable (press \texttt{A} for axes).

\subsection{Step-by-Step Mode}

The rules have sequential structure that autonomous animation loses.
For example, Kan filling has three stages: (1) show the open box, (2) show the
given data $[\varphi \mapsto u]$, (3) fill the lid. A ``step'' button would let
users pause at each conceptual stage, with a text callout explaining what just
happened.

\subsection{Proof Relevance: Show the Witness}

For rules like Composition and Transport, the \emph{proof term} (the witness) is as
important as the judgment. Consider rendering the term in a ``proof tree'' panel
alongside the 3D view:

\begin{itemize}
  \item Split the content area: left = 3D scene (60\%), right = proof tree (40\%).
  \item The proof tree uses KaTeX and collapses/expands nodes.
  \item Hovering a node in the proof tree highlights the corresponding object in
    3D; the existing hover system already supports this pattern.
\end{itemize}

% ────────────────────────────────────────────────────────────────────────────
\section{New Concepts to Add}
% ────────────────────────────────────────────────────────────────────────────

\subsection{Path Concatenation and Inversion}

Before composition (which is the general $n$-dimensional operation), show the
simpler 1D cases:
\begin{itemize}
  \item \textbf{Concatenation}: $p \cdot q$ where $p : a = b$ and $q : b = c$.
    Show as two line segments being joined end-to-end.
  \item \textbf{Inversion}: $p^{-1} : b = a$.
    Show as the path being ``reversed'' (arrow flips direction).
  \item \textbf{Reflexivity}: $\mathrm{refl}_a : a = a$.
    Show as a degenerate constant path.
\end{itemize}

\subsection{Higher-Dimensional Paths}

The square and cube primitives exist but are not connected to their type-theoretic
meaning. Add pages for:
\begin{itemize}
  \item \textbf{Square}: A term $p : \mathrm{Path}\,(\mathrm{Path}\,A\,a\,b)\,q\,r$
    as a square face with all four edges labeled.
  \item \textbf{Cube}: A 3-dimensional path as a solid with labeled faces showing
    the coherence conditions.
\end{itemize}

\subsection{The Face Lattice $\mathbb{F}$}

The face formula type $\mathbb{F}$ is used in Composition, Kan filling, and Glue
types but never visualized. A dedicated page showing:
\begin{itemize}
  \item The Hasse diagram of $\mathbb{F}$ as a lattice.
  \item Interactive toggle: click a face of a cube and see the corresponding formula.
  \item The meet/join operations $\varphi \land \psi$, $\varphi \lor \psi$.
\end{itemize}

\subsection{Univalence Proof via Glue}

The classical proof of univalence in cubical type theory uses Glue types. Walking
through this proof interactively would be a flagship feature:
\begin{enumerate}
  \item Start with an equivalence $e : A \simeq B$.
  \item Construct $\mathrm{Glue}\,[i=0 \mapsto (A, e)]\,B$.
  \item Show that the resulting type gives a path $A = B$ in $\mathcal{U}$.
  \item Animate the passage through the four corners of a square in the universe.
\end{enumerate}

% ────────────────────────────────────────────────────────────────────────────
\section{UX and Navigation}
% ────────────────────────────────────────────────────────────────────────────

\subsection{Dependency Graph View}

The sidebar currently lists rules as a flat categorized list. Adding a
\textbf{dependency graph view} (toggle-able) would show how concepts build on each
other:
\[
  \mathbb{I} \;\longrightarrow\; \text{Path Type} \;\longrightarrow\; \text{Transport} \;\longrightarrow\; \text{Composition} \;\longrightarrow\; \text{Kan Filling} \;\longrightarrow\; \text{Glue Types}
\]
Clicking a node in the graph navigates to that rule. Edges should be labeled with
the dependency relationship.

\subsection{URL-Based Deep Linking}

Currently the visualizer always starts at the Interval page. Adding URL hash routing
(e.g.\ \texttt{/cubical-viz/\#transport-rule}) would allow linking directly to a
specific rule, making it useful for sharing in papers or lecture notes.

\subsection{Gradient Colors for Continuous Transformations}

Many of the concepts in cubical type theory involve \emph{continuous variation} along
the interval $\mathbb{I}$. Flat colors fail to communicate this — they make
transformations look like discrete jumps. Vertex-color gradients can make the
continuity viscerally obvious.

\begin{itemize}
  \item \textbf{Interval dimension}: Color geometry by the value of $i \in [0,1]$
    using a gradient from cyan ($i=0$, \texttt{\#00DDFF}) to orange ($i=1$,
    \texttt{\#FF8800}). Any object parameterized by $i$ should inherit this gradient.
  \item \textbf{Transport path}: The arc from $A(i_0)$ to $A(i_1)$ should be
    gradient-colored along its length, making it clear which end is the source and
    which is the target.
  \item \textbf{Composition square}: The square face being filled should use a 2D
    gradient: $i$-axis (horizontal) and $j$-axis (vertical) each contribute a
    color channel, so the entire face reads as a color map of $(i, j) \in [0,1]^2$.
  \item \textbf{Morphing fibers in Transport}: As the type family $A$ deforms, the
    geometry should be colored with a gradient that tracks \emph{how far along the
    path} we are — the color encodes $i$ continuously so the viewer can see the
    transition is smooth, not a teleport.
  \item \textbf{Kan filling lid}: The lid being filled should be colored with a
    radial gradient emanating from the edges (where data is given) toward the
    center (where it is being computed), suggesting that the fill ``spreads inward''
    from known boundary data.
\end{itemize}

In Three.js this is achieved with \texttt{BufferGeometry} vertex colors
(\texttt{geometry.setAttribute('color', ...)}) combined with
\texttt{vertexColors: true} on \texttt{MeshBasicMaterial}. For line objects,
\texttt{LineBasicMaterial} supports vertex colors equally well.

\subsection{Camera Presets}

For each rule, define 2--3 \emph{camera presets} (top view, front view, isometric)
that the user can jump to with keyboard shortcuts (1, 2, 3). This helps when the
auto-orbit leaves the scene at an unhelpful angle.

% ────────────────────────────────────────────────────────────────────────────
\section{Priority Summary}
% ────────────────────────────────────────────────────────────────────────────

\begin{center}
\begin{tabular}{llll}
\toprule
\textbf{Item} & \textbf{Impact} & \textbf{Effort} & \textbf{Priority} \\
\midrule
Color consistency across rules & Medium & Low & \textbf{High} \\
Gradient colors for continuous transforms & High & Low & \textbf{High} \\
Axis labels with dimension variables & Medium & Low & \textbf{High} \\
URL hash routing & Low & Low & \textbf{High} \\
Composition as square fill & High & Medium & \textbf{High} \\
Transport with morphing fibers & High & Medium & \textbf{High} \\
Step-by-step mode & High & Medium & Medium \\
Glue types: bidirectional arrows & Medium & Medium & Medium \\
Kan filling: face formula $\varphi$ visual & Medium & Medium & Medium \\
Proof tree panel & High & High & Medium \\
Path concatenation / inversion pages & Medium & Medium & Medium \\
Face lattice $\mathbb{F}$ page & High & High & Low \\
Univalence proof walkthrough & Very High & Very High & Low \\
Dependency graph navigation & Medium & High & Low \\
\bottomrule
\end{tabular}
\end{center}

% ────────────────────────────────────────────────────────────────────────────
\section{Conclusion}
% ────────────────────────────────────────────────────────────────────────────

The most impactful short-term improvements are:
\begin{enumerate}
  \item Fixing Composition to show an actual square being filled (not just an arc).
  \item Making Transport show the \emph{varying type family} as morphing geometry.
  \item Establishing a consistent color vocabulary so users can read across rules.
  \item Adding axis labels so spatial dimensions map clearly to interval variables.
\end{enumerate}

These four changes would make the tool substantially more informative for someone
learning cubical type theory. The longer-term additions (univalence proof walkthrough,
face lattice, proof trees) would make it genuinely useful as a pedagogical companion
to papers like CCHM (Cohen--Coquand--Huber--M\"ortberg) or the HoTT book's cubical
chapter.

\end{document}
